% !TEX root = main21.tex
\section{Wider-grid calibration and VIX implementation}
\label{app:wide-grid-details}

\subsection{State representation and checks}
The wider model has 200 interior strikes plus two boundary nodes; the
boundary levels remain one half of the smallest interior strike and one
and a half times the largest. All 25 expiries share that support on a given
historical date. The support floats between historical dates but cash
strikes remain fixed within a simulated path and within each historical
one-session alignment. The DLV-squared parameterization retains minimum
$10^{-4}$ and upper cap $9$, with the original log-space barrier. The
completed fit's largest DLV-squared value is $6.319$, with no coordinate
at the cap. Across all historical states, independent float64 reconstruction
has minimum probability $3.46\times10^{-24}$ and maximum mass/forward errors
below $4.3\times10^{-14}$.

Training uses float32 as in the baseline. Its finite-difference point-mass
initializer has a small signed boundary roundoff of about
$-1.79\times10^{-7}$. Such checkpoints are accepted only when negative
mass is bounded by that initializer's roundoff and independent float64
reconstruction passes strict positivity, unit-mass and unit-forward checks.
The Q simulation and its initial-state reconstruction use float64.
This numerical calibration check is separate from the dynamic drift tests.

\subsection{Diffusion and audit conventions}
\label{sec:wide-audit-definition}
Historical log-DLV changes are aligned at fixed expiry and cash strike;
the endpoint surface is interpolated at origin relative strike divided by
the realized gross spot return. A nodewise physical AR fit removes the
conditional mean before smoothing and the joint spot/surface PCA.
For the Q audit the drift is recalculated from the density recursion,
not obtained by setting the empirical physical drift to zero.
The simulator retains the paper's log-Euler step and its exceptional
relative-drift taming threshold $0.5$; it is not an exact finite-step
martingale correction. The independent 5\% Black factor has zero Q drift.
The two conditional objectives, fold construction and spot normalization
are those of the main paper; each
fold uses a 30-epoch affine warm start and 30-epoch affine-plus-deep fit,
with two width-32 SiLU layers, learning rate $10^{-3}$ and batch size 2,048.
The simulation seed is 20261024. Constant physical risk prices are also
re-estimated, but the new reported path diagnostics are under Q only.

For each live day--instrument coordinate, the return is
$R_{p,n}^i=(V_{p,n+1}^i-V_{p,n}^i)/Z_{p,n}$, where $Z$ is pre-step spot
and $V^i$ is spot or an option price. The unconditional score is its sample
mean divided by the standard error computed across 50,000 antithetic-pair
averages. Table~\ref{tab:wide-q-audit} reports the maximum and empirical
99th percentile of $|z|$ over the 570 coordinates, together with counts
$N_c$ for $|z|>c$. These coordinatewise thresholds are not family-wise
$p$-values.

Conditional policies use current Brownian-factor histories, log spot,
live log option-price-to-spot ratios, time and remaining maturities. For
an action $a(X)$, write $r=a(X)^\top R$. The squared objective maximizes
the training mean of $r-\tfrac12\lVert a\rVert^2-\tfrac12r^2$; the entropy
objective minimizes $\log\operatorname{mean}\exp(-r+\tfrac12\lVert a\rVert^2)$.
Complete antithetic pairs are split into two equal folds; input
standardization and action fitting use only the training fold, and each
frozen policy is evaluated on the other fold. The folds are then exchanged.
For either objective, the reported $T$ is the mean out-of-fold daily
portfolio return divided by its pair-based standard error, after averaging
over the 20 days within each path. The main paper's per-test colour rule is
black for $|T|\leq1.96$, orange for $1.96<|T|\leq2.576$ and red above
$2.576$. Both new conditional statistics are therefore black.

\subsection{Market-selected discrete variance}
For each eligible expiry let $D_T$ be the discount factor, $F_T$ the
put--call-parity forward, $K_0$ the greatest available strike not exceeding
$F_T$, and $O_T(K)$ the selected OTM midpoint price (the average of call
and put prices at $K_0$). The Cboe total-variance formula is
\begin{equation}\label{eq:wide-vix-strip}
  v_T^{\mathrm{strip}}
  =\frac{2}{D_T}\sum_i\frac{\Delta K_i}{K_i^2}O_T(K_i)
   -\left(\frac{F_T}{K_0}-1\right)^2.
\end{equation}
Interior $\Delta K_i$ are half the neighbouring-strike difference; endpoints
use the adjacent-strike difference. Selection walks outward from $K_0$,
omits zero-bid contracts and stops after two consecutive zero bids on each
side. Friday expiries strictly between 23 and 37 calendar days are selected
with their AM/PM settlement convention. For minute counts $N_1<N_{30}<N_2$,
\begin{equation}\label{eq:wide-vix-interpolation}
  \mathrm{VIX}
  =100\sqrt{\frac{N_{365}}{N_{30}}\left[
    \frac{N_2-N_{30}}{N_2-N_1}v_{T_1}
    +\frac{N_{30}-N_1}{N_2-N_1}v_{T_2}\right]},
\end{equation}
where $N_{30}=43,200$ and $N_{365}=525,600$.
These conventions follow the \href{https://cdn.cboe.com/resources/vix/VIX_Methodology.pdf}
{Cboe VIX methodology}, pp.~2--7.

The model replaces only the option prices in
\eqref{eq:wide-vix-strip}; the observed raw bid cutoffs, forward, $K_0$,
discount factors, contracts and interpolation times are shared by both
GROMM variants. No fictitious model bid/ask quotes are introduced.
Filtered-market VIX uses the synthetic call-equivalent bid/ask midpoints
passed to SANOS, including parity conversion for puts, not fitted option
prices. Since preprocessing has removed zero bids, the filtered-strip
calculation respects the observed raw-book cutoffs, then computes its own
forward, $K_0$ and strike quadrature on the surviving input quotes.

\subsection{Analytic full-strike variance}
For the forward-normalized mixture $S_T/F_T=s_JB_T$, with probabilities
$p_j(T)$ on positive support $s_j$ and independent mean-one Black factor
of business-time variance $0.05^2\tau_b(T)$, the model gives
\begin{equation}\label{eq:wide-analytic-variance}
  v_T^{\mathrm{analytic}}
  =-2\E^\Q\!\left[\log(S_T/F_T)\right]
  =-2\sum_jp_j(T)\log s_j+0.05^2\tau_b(T).
\end{equation}
This is the full-strike log-contract/variance-swap quantity, not a discrete
VIX strip. We substitute these two expiry values into
\eqref{eq:wide-vix-interpolation}; a direct 30-calendar-day model evaluation
has a different maturity convention and is not used in
Table~\ref{tab:wide-vix} or Figure~\ref{fig:wide-vix-comparison}.

Calendar time is measured from 16:15 New York to 09:30 for AM settlement
or 16:00 for PM settlement. Model business time uses the NYSE session
calendar divided by 252 and the implementation's $0.8$-day AM subtraction.
Fractional model maturities interpolate adjacent daily densities. Discount
factors use the cached OptionMetrics zero curve with linear rate
interpolation, not Cboe's Treasury spline. No maturity beyond 25 business
days is extrapolated. These clock and input differences are retained in
the reported comparisons, not calibrated away by a VIX rescaling.

\subsection{Reproducibility}
The fresh calibration is run by
\path{quantdata/sdlv/daily/wide_fit.py}; the Q audit by
\path{quantdata/sdlv/daily/wide_audit.py}; and the VIX calculations by
\path{quantdata/sdlv/daily/vix_comparison.py} and
\path{quantdata/dysanos/filtered_vix.py}, using the shared formula in
\path{quantdata/vixdata/formula.py}.
The source for the new tables and figure is
\path{quantdata/sdlv/daily/paper_wide_grid.py}.
The retained \path{generated/wide_grid/metrics.json} records the completed
fit and audit, source paths and SHA-256 hashes; the accompanying
\path{vix_series.csv} supplies the plotted observations. Large training
data, checkpoints and simulation chunks remain outside the Git tree.
This experiment uses training run \path{20261023-kVyePV} and Q-audit
run \path{q-audit-20261024-4mkznc}. The VIX and Q results are checked
to refer to the identical final 1,000-epoch checkpoint; no interim VIX
snapshot is used in this addendum.
